\reading{CR-3}{Credit Risk Management}{DEFER}{1}

\begin{crkeybox}[Learning objectives for this reading]
\begin{enumerate}[label=\textbf{\alph*.},leftmargin=2em]
\item Describe key elements of an effective lending or financing policy.
\item Explain the importance and challenges of setting exposure and concentration limits.
\item Describe the scope and allocation processes of a bank's credit facility and explain bank-specific policies and actions to reduce credit risk.
\item Discuss factors that should be considered during the credit asset classification process.
\item Describe and explain loan loss provisions and loan loss reserves.
\item Identify and explain the components of expected loss and differentiate between expected loss and unexpected loss.
\item Explain the requirements for estimating expected loss under IFRS 9.
\item Describe a workout procedure for loss assets and compare the following two approaches used to manage loss assets: retaining loss assets and writing off loss assets.
\item Explain the components of credit risk analysis.
\item Explain the components of credit risk management capacity, and identify key questions that the board of directors of a bank should ask.
\end{enumerate}
\end{crkeybox}

% =====================================================================
\lo{CR-3 a}{Describe key elements of an effective lending or financing policy.}

\begin{crdefbox}[Credit risk, and why the lending policy matters]
Credit risk (i.e., counterparty risk) is the chance that a borrower will either
make delayed payments or no payments at all. Credit risk is linked to
approximately \term{70\% of a bank's balance sheet}, which makes credit risk the
primary cause of bank failure.
\end{crdefbox}

An effective lending or financing policy is crucial for managing credit risk,
which can lead to delayed or non-payments by borrowers and impact a bank's
liquidity, contributing to bank failures.

Key elements include:

\begin{itemize}
\item defining credit facilities and risk management processes;
\item emphasizing flexibility;
\item complying with regulations;
\item classifying assets by risk exposure; and
\item establishing provisions for potential losses.
\end{itemize}

This policy ensures proactive risk management, regulatory compliance, and
financial stability.

% =====================================================================
\lo{CR-3 b}{Explain the importance and challenges of setting exposure and concentration limits.}

Regulators focus on limiting credit risk exposures across three key areas:
exposure to individual customers, related-party financing, and overexposure to
specific economic sectors or regions.

\begin{enumerate}
\item \term{Exposure.} The maximum loan amount to a single borrower (e.g.,
10--25\% of capital).
\item \term{Concentration.} The maximum exposure to a sector or region.
\item \term{Related parties.} The maximum loan amount to affiliated companies or
shareholders.
\end{enumerate}

\begin{itemize}
\item \term{Concentration limits} dictate the maximum exposure to a single
client or sector, typically set as a percentage of the bank's reserves or
capital, and prevent overreliance on a single client, sector, or region.
Challenges arise with less-direct forms of credit, requiring careful monitoring
to avoid breaching these limits.
\item \term{Vigilance in related-party financing} is crucial to prevent
conflicts of interest. Regulators often impose limits as a percentage of Tier 1
capital, or banks establish internal policies for oversight.
\end{itemize}

Overall, robust policies and regulatory oversight are vital for managing credit
risk and ensuring the stability of financial institutions. Concentration limits
and monitoring of related-party transactions are key components of prudent
lending practices.

% =====================================================================
\lo{CR-3 c}{Describe the scope and allocation processes of a bank's credit facility and explain bank-specific policies and actions to reduce credit risk.}

The credit facility management of a bank involves the entire process of loan
origination, appraisal, supervision, and collection, with the necessary
flexibility to adapt to changing market conditions.

Key factors in lending policy and risk reduction include:

\begin{enumerate}
\item \term{Lending authority.} Establishing limits for lending officers,
possibly with higher limits for those with more experience or tenure. This
process may be centralized for smaller banks and decentralized for larger banks
based on factors like geographical area or customer types.
\item \term{Types and distribution of loans.} Defining which credit instruments
the bank will offer based on lending officers' expertise, deposit structure, and
credit demand. Limitations are typically imposed for lending categories, with
deviations requiring board approval.
\item \term{Appraisal process.} Formally defining the appraisal process,
including limits on appraisals by credit type, reappraisal processes for renewals
and extensions, and acceptable loan-to-value ratios.
\item \term{Loan pricing.} Structuring loan rates by type, considering cost of
funds, supervision costs, and probability of loss. Periodic review of rates to
adjust to market conditions is essential.
\item \term{Maturities.} Establishing clear maximum maturities for various loan
types, considering repayment sources and collateral useful life.
\item \term{Exposure concentration.} Monitoring risks associated with geographic
or sector-level concentration, including international lending, to understand and
mitigate unique risk factors.
\item \term{Availability of financial information.} Specifying required financial
information for commercial and consumer loans, involving external credit checks,
and aligning financial projections with loan maturity. Mandatory for assessing
creditworthiness and loan suitability.
\item \term{Monitoring of collections.} Addressing processes for dealing with
delinquencies, including reporting to the board and determining potential losses
and alternative repayment plans.
\item \term{Limits on total loans outstanding.} Setting policy limits on total
loans outstanding relative to capital, deposits, or total assets, factoring in
deposit volatility and credit risks to reduce overall risk exposure.
\item \term{Maximum loan-to-value ratio.} Establishing margin requirements for
all security types and including provisions for periodic repricing of collateral
to reduce risk.
\item \term{Impairment recognition.} Periodically and systematically identifying
loan impairments to address loans deemed uncollectible.
\item \term{Renegotiated debt treatment.} Managing loans restructured to reduce
interest rates or principal amounts, monitoring restructurings as an indication
of serious credit risk issues.
\item \term{Written internal guidelines.} Proper documentation of all lending
policies to ensure clear understanding and consistent implementation across the
organization, and for regulatory compliance.
\end{enumerate}

% =====================================================================
\lo{CR-3 d}{Discuss factors that should be considered during the credit asset classification process.}

Regulators typically classify assets into five categories.

\begin{enumerate}
\item \term{Standard (pass).} Assets labeled as standard are not delinquent,
often fully secured with cash equivalents regardless of other credit factors.
Loans secure and without delinquency concerns.
\item \term{Specially mentioned (watch).} Assets with potential weaknesses
impacting repayment capacity fall into this category, with varied reasons for
expected weakness.
\item \term{Substandard.} Loans exhibiting defined credit weaknesses
jeopardizing repayment, often with primary repayment sources failing and the bank
relying on collateral. Nonperforming loans over 90 days delinquent typically fall
here.
\item \term{Doubtful.} Similar to substandard assets but with a higher
expectation of loss. Nonperforming loans over 180 days delinquent, unless well
secured, are categorized as doubtful, though some hope for repayment remains.
\item \term{Loss.} Assets where delinquency is almost certain, prompting
reclassification as a loss. This does not imply zero chance of collection but
indicates the impracticality of delaying write-off. Nonperforming loans over one
year delinquent typically fall into this category, unless adequately secured.
\end{enumerate}

\begin{center}
\begin{tikzpicture}[font=\sffamily\footnotesize]
\def\w{2.92}
\foreach \i/\lab/\col/\sub/\thr in {
  0/{Standard\\(pass)}/FRMGreen/{not delinquent}/{no day threshold},
  1/{Specially\\mentioned}/FRMGreen/{potential weakness}/{no day threshold},
  2/{Substandard}/FRMGold/{defined weakness}/{over 90 days},
  3/{Doubtful}/FRMRed/{higher loss expected}/{over 180 days},
  4/{Loss}/FRMRed/{write-off imminent}/{over 1 year}}
{
  \node[draw=\col, fill=\col!8, line width=0.8pt, rounded corners=1pt,
        text width=2.55cm, align=center, minimum height=1.3cm]
        (c\i) at ({\i*\w},0) {\scriptsize\textbf{\lab}\\[1pt]\scriptsize\sub};
  \node[draw=FRMRule, fill=white, line width=0.6pt, text width=2.55cm,
        align=center, minimum height=0.72cm] (t\i) at ({\i*\w},-1.55)
        {\scriptsize\thr};
  \draw[FRMRule, line width=0.6pt] (c\i.south) -- (t\i.north);
}
\node[left, font=\scriptsize\bfseries, color=FRMNavy] at (-1.5,-1.55)
  {Delinquency};
\draw[-{Stealth[length=5pt]}, FRMGrey, line width=0.8pt]
  ($(c0.north)+(0,0.35)$) -- ($(c4.north)+(0,0.35)$)
  node[midway, above, font=\scriptsize\itshape, text=FRMGrey, fill=white,
       inner sep=1.5pt] {increasing severity};
\end{tikzpicture}
\end{center}
\figcap{The classification ladder and the delinquency threshold attaching to each
rung. The thresholds begin only at substandard: the first two categories are
judgements about repayment capacity, not about elapsed days. Each threshold also
carries an escape clause --- doubtful applies unless the loan is well secured,
and loss applies unless it is adequately secured.}

\begin{crkeybox}[Reviewing non-performing loans]
The review considers the age of the delinquency, the reasons for the
delinquency, a case-by-case assessment, and the setting of provision levels.
\end{crkeybox}

% =====================================================================
\lo{CR-3 e}{Describe and explain loan loss provisions and loan loss reserves.}

\begin{crdefbox}[Provisions versus reserves]
\term{Loan loss provisions} are an expense set aside each year to cover potential
future losses on loans.

\smallskip
\term{Loan loss reserves} are the accumulated amount of loan loss provisions over
time, recorded on the bank's balance sheet.
\end{crdefbox}

\textbf{Factors affecting reserves:}

\begin{center}
\footnotesize
\begin{tabular}{L{7.2cm} L{7.2cm}}
\toprule
Credit quality policies and procedures & Previous loss exposures \\
Growth in the loan book & Quality of management in charge of lending \\
Loan collection and recovery policies & Changes in the economic environment and
business conditions \\
\bottomrule
\end{tabular}
\end{center}

% =====================================================================
\lo{CR-3 f}{Identify and explain the components of expected loss and differentiate between expected loss and unexpected loss.}

In the context of loan loss scenarios, \term{expected loss (EL)} is the
anticipated dollar loss if a borrower defaults. It comprises three components.

\begin{enumerate}
\item \term{Probability of default (PD).} The likelihood of a borrower failing to
make timely payments. This measure considers \emph{who} the bank has lent money
to.
\item \term{Loss given default (LGD).} The expected percentage loss of the amount
borrowed in the event of default. This measure considers \emph{what product} was
used for the credit risk.
\item \term{Exposure at default (EAD).} An estimate of the dollar loss exposure
of a bank when a borrower defaults.
\end{enumerate}

\begin{crfmlbox}[Expected loss]
\[ \mathrm{EL}(\$) \;=\; \mathrm{PD}(\%) \times \mathrm{LGD}(\%) \times \mathrm{EAD}(\$) \]
\vk{$\mathrm{PD}$ = probability of default, a percentage;
$\mathrm{LGD}$ = loss given default, a percentage of the amount borrowed;
$\mathrm{EAD}$ = exposure at default, a currency amount. The product of two
percentages and one currency amount returns a currency amount.}
\end{crfmlbox}

\begin{center}
\footnotesize
\begin{tabular}{c c r r r r r}
\toprule
\thead{Borrower} & \thead{Rating} & \thead{Loan} & \thead{PD} & \thead{LGD} &
\thead{EAD} & \thead{EL} \\
\midrule
1 & AAA & \$1,000,000 & 0.09\% & 30.0\% & \$550,000 & \$148.50 \\
2 & B+ & \$1,000,000 & 3.78\% & 45.0\% & \$785,000 & \$13,352.85 \\
3 & D (defaulted) & \$1,000,000 & 100.00\% & 100.0\% & \$1,000,000 & \$1,000,000 \\
\bottomrule
\end{tabular}
\end{center}
\figcap{Expected loss components. Note that the loan amount plays no part in the
arithmetic: all three loans are \$1,000,000, but EL is driven by EAD, which
differs from the loan in the first two rows and equals it only on default.}

\begin{crdefbox}[Unexpected loss]
In a very simplistic way, an \term{unexpected loss} is a loss that was not
captured by the expected loss formula. This type of loss is typically thought of
as a tail risk event (i.e., a 99th percentile event).
\end{crdefbox}

\begin{center}
\begin{tikzpicture}
\begin{axis}[
  width=13.2cm, height=5.6cm,
  axis lines=left, xlabel={Loss over the horizon}, ylabel={Probability},
  xlabel style={font=\sffamily\scriptsize}, ylabel style={font=\sffamily\scriptsize},
  xtick=\empty, ytick=\empty,
  xmin=0, xmax=7, ymin=0, ymax=0.30,
  clip=true]
\addplot[domain=0:7, samples=200, smooth, draw=FRMNavy, line width=1pt,
  name path=curve] {0.55*x^1.4*exp(-1.15*x)};
\path[name path=base] (axis cs:0,0) -- (axis cs:7,0);
\addplot[FRMRed!35] fill between[of=curve and base, soft clip={domain=4.0:7}];
% EL line
\draw[FRMGreen, line width=0.9pt] (axis cs:2.09,0) -- (axis cs:2.09,0.22);
\node[above, font=\sffamily\scriptsize\bfseries, text=FRMGreen, fill=white,
  inner sep=1.5pt] at (axis cs:2.09,0.22) {Expected loss};
% percentile line
\draw[FRMRed, line width=0.9pt] (axis cs:4.0,0) -- (axis cs:4.0,0.15);
\node[above, font=\sffamily\scriptsize\bfseries, text=FRMRed, fill=white,
  inner sep=1.5pt] at (axis cs:4.55,0.15) {99th percentile};
\node[font=\sffamily\scriptsize, text=FRMRed, fill=white, inner sep=1.5pt]
  at (axis cs:5.55,0.085) {unexpected loss};
\draw[-{Stealth[length=4pt]}, FRMRed] (axis cs:5.2,0.07) -- (axis cs:4.7,0.022);
\end{axis}
\end{tikzpicture}
\end{center}
\figcap{Illustrative shape only, no data. Expected loss sits near the body of the
distribution and is covered by provisioning; unexpected loss is what lives out in
the tail beyond it. The two are not two sizes of the same thing --- EL is a cost
you price for, UL is a risk you hold capital against.}

% =====================================================================
\lo{CR-3 g}{Explain the requirements for estimating expected loss under IFRS 9.}

Before 2018, International Accounting Standard (IAS) 39 required an
\term{incurred loss model}. Beginning 1 January 2018, International Financial
Reporting Standards (IFRS) 9 switched the requirement to a three-stage
\term{expected loss model}.

\begin{itemize}
\item \term{Stage 1 --- performing assets.} All performing assets (i.e., no
delinquency) should have provisions for expected loss calculated using a
12-month expected loss methodology. This analysis should include effective
interest based on the \emph{gross} amount.
\item \term{Stage 2 --- delinquent assets.} Assets with any level of delinquency
should use a \emph{lifetime} expected loss model, including effective interest
based on the \emph{gross} amount.
\item \term{Stage 3 --- nonperforming assets.} Nonperforming assets should also
use a \emph{lifetime} expected loss model, including effective interest based on
the \emph{net} (carrying) amount.
\end{itemize}

\begin{center}
\begin{tikzpicture}[font=\sffamily\footnotesize]
\def\w{4.3}
\foreach \i/\lab/\st in {0/{Stage 1}/{performing,\\no delinquency},
                         1/{Stage 2}/{any level of\\delinquency},
                         2/{Stage 3}/{nonperforming}}{
  \node[draw=FRMNavy, fill=PaleNavy, line width=0.7pt, rounded corners=1pt,
        text width=3.9cm, align=center, minimum height=1.0cm]
        at ({\i*\w},0) {\textbf{\lab}\\[1pt]\scriptsize\st};
}
\node[left, font=\scriptsize\bfseries, color=FRMNavy] at (-2.05,-1.5)
  {ECL horizon};
\node[draw=FRMGreen, fill=PaleGreen, line width=0.7pt, text width=3.9cm,
  align=center, minimum height=0.85cm] at (0,-1.5) {\scriptsize 12-month};
\node[draw=FRMGold, fill=PaleGold, line width=0.7pt, text width=8.2cm,
  align=center, minimum height=0.85cm] at (6.45,-1.5) {\scriptsize lifetime};
\node[left, font=\scriptsize\bfseries, color=FRMNavy] at (-2.05,-2.7)
  {Interest base};
\node[draw=FRMGreen, fill=PaleGreen, line width=0.7pt, text width=8.2cm,
  align=center, minimum height=0.85cm] at (2.15,-2.7) {\scriptsize gross amount};
\node[draw=FRMRed, fill=PaleRed, line width=0.7pt, text width=3.9cm,
  align=center, minimum height=0.85cm] at (8.6,-2.7) {\scriptsize net (carrying) amount};
\end{tikzpicture}
\end{center}
\figcap{The two attributes switch at \emph{different} boundaries. The expected
loss horizon changes between stage 1 and stage 2; the effective interest base
changes between stage 2 and stage 3. Stage 2 is therefore the only stage that
pairs a lifetime horizon with a gross interest base, which is exactly the cell an
examiner can swap without the sentence sounding wrong.}

% =====================================================================
\lo{CR-3 h}{Describe a workout procedure for loss assets and compare the following two approaches used to manage loss assets: retaining loss assets and writing off loss assets.}

The loan workout procedure is a very important part of reducing risk for a bank.
It is a strategic approach banks use to mitigate risk by rehabilitating loans
that are in jeopardy of default. It is helpful to review historical attempts to
collect and the success rate.

A common workout procedure may include the following four steps, which can be
applied \emph{in any order}:

\begin{enumerate}
\item \term{Reduce the bank's credit risk exposure} by collecting additional
capital, collateral, or guarantees from the borrower.
\item \term{Work with the borrower} to identify areas for potential improvement
(e.g., providing advice, developing a plan to reduce costs and increase earnings,
selling assets, or restructuring the loan). This option borders on consulting.
\item \term{Introduce a third party} as a possible joint-venture partner or a
takeover --- find a partner or buyer for the borrower's business.
\item \term{Liquidation} through out-of-court settlement. This could involve
foreclosure, liquidating collateral, or calling on guarantees.
\end{enumerate}

\subsection*{Handling loss assets: two traditions}

\begin{center}
\begin{tikzpicture}[font=\sffamily\footnotesize]
\node[draw=FRMNavy, fill=PaleNavy, line width=0.8pt, rounded corners=1pt,
      text width=6.4cm, align=left, inner sep=7pt] (ret) at (0,0)
  {\textbf{RETENTION}\ \ \scriptsize(British tradition)\\[3pt]
   \scriptsize Keep loss assets on the books to allow time for recovery efforts.};
\node[draw=FRMGold, fill=PaleGold, line width=0.8pt, rounded corners=1pt,
      text width=6.4cm, align=left, inner sep=7pt] (wo) at (8.4,0)
  {\textbf{WRITING OFF}\ \ \scriptsize(U.S. model)\\[3pt]
   \scriptsize Immediately remove losses from the balance sheet.};
\node[draw=FRMRule, fill=white, line width=0.6pt, text width=6.4cm,
      align=center, inner sep=6pt] (reteff) at (0,-2.05)
  {\scriptsize effect on the balance sheet\\[2pt]
   \textbf{\color{FRMGreen}loss reserve increases}};
\node[draw=FRMRule, fill=white, line width=0.6pt, text width=6.4cm,
      align=center, inner sep=6pt] (woeff) at (8.4,-2.05)
  {\scriptsize effect on the balance sheet\\[2pt]
   \textbf{\color{FRMRed}apparent size of loan loss reserves falls}};
\draw[-{Stealth[length=5pt]}, FRMGreen, line width=0.9pt] (ret) -- (reteff);
\draw[-{Stealth[length=5pt]}, FRMRed, line width=0.9pt] (wo) -- (woeff);
\end{tikzpicture}
\end{center}
\figcap{The same underlying loss, two opposite reported outcomes. Retaining the
asset raises the loss reserve on the balance sheet; writing it off reduces the
apparent size of the loan loss reserves. The direction of the reserve movement is
the discriminating fact.}

% =====================================================================
\lo{CR-3 i}{Explain the components of credit risk analysis.}

Credit risk analysis should evaluate which lending products are being used, who
the customers are, and for how long money is being lent. It analyzes lending
products, customers, and loan terms to assess risk.

Credit risk analysis involves:

\begin{enumerate}
\item Summarizing major loan types, including customer numbers, maturity, and
interest rates.
\item Detailing the loan portfolio distribution by factors like number of loans,
amounts, currencies, maturities, industries, and borrower category.
\item Listing loans with guarantees from governmental bodies.
\item Reviewing loans based on risk classification.
\item Analyzing nonperforming loans by vintage year for trend identification.
\end{enumerate}

% =====================================================================
\lo{CR-3 j}{Explain the components of credit risk management capacity, and identify key questions that the board of directors of a bank should ask.}

Credit risk management capacity reviews should consider the lending process, the
bank's staffing, and its flow of information. The end goal of information flow
analysis is to make sure that senior managers and the Board have the information
they need in a timely manner.

\subsection*{1) Lending process}
\begin{enumerate}[label=\alph*)]
\item Are lending policies well-organized and clearly defined?
\item What is the ratio of credit applications appraised to those approved, and
how does it reflect risk assessment practices?
\item How are exceptions to lending policies handled, and is there a consistent
process for addressing them?
\end{enumerate}

\subsection*{2) Staffing}
\begin{enumerate}[label=\alph*)]
\item Who is involved in various stages of the lending process, including
origination, appraisal, supervision, and monitoring?
\item How many staff are dedicated to these roles, what are their
qualifications, and how frequently are they trained to stay updated on risk
management practices?
\item Is there a sufficient number of skilled personnel to handle the volume and
complexity of lending activities?
\end{enumerate}

\subsection*{3) Information flows}
\begin{enumerate}[label=\alph*)]
\item How is lending-related information disseminated within the organization,
particularly to senior management and the Board?
\item Is the flow of information timely, ensuring that decision-makers have
access to relevant data when assessing credit risk?
\item Are mechanisms in place to ensure the accuracy and reliability of the
information provided, and is it cost-effective to maintain these systems?
\end{enumerate}
